\chapter{The reacting Navier--Stokes equations}
\label{ch:governing}

The onboarding introduction lists four conservation laws. Here each one is built from the generic
balance \eqref{eq:generic}, so you can see where every term comes from. We then simplify them
for the low-speed flames Cantera solves.

\section{Mass}

With $\varphi=1$, no diffusion of total mass and no source:
\begin{equation}
  \pd{\rho}{t}+\divg(\rho\vu)=0 .
  \label{eq:mass}
\end{equation}

\begin{symbols}
$\rho$ & density & \si{\kilo\gram\per\metre\cubed}\\
$\vu$ & velocity & \si{\metre\per\second}\\
$t$ & time & \si{\second}\\
\end{symbols}

\section{Species}

With $\varphi=Y_k$, diffusive flux $\vj_k=\rho Y_k\vV_k$ and chemical source $W_k\wdot_k$:
\begin{equation}
  \pd{(\rho Y_k)}{t}+\divg(\rho\vu Y_k)=-\divg(\rho Y_k\vV_k)+W_k\wdot_k,\qquad k=1\dots K.
  \label{eq:species}
\end{equation}

\begin{symbols}
$Y_k$ & mass fraction of species $k$ & \nodim\\
$\vV_k$ & diffusion velocity of $k$ & \si{\metre\per\second}\\
$\rho Y_k\vV_k$ & diffusive mass flux of $k$ & \si{\kilo\gram\per\metre\squared\per\second}\\
$W_k\wdot_k$ & chemical mass production rate of $k$ & \si{\kilo\gram\per\metre\cubed\per\second}\\
\end{symbols}
Summing over all $k$, and using $\sum Y_k=1$, $\sum\rho Y_k\vV_k=0$ and $\sum W_k\wdot_k=0$, gives back
\eqref{eq:mass}. So only $K-1$ species equations are independent. Cantera solves all $K$ and
enforces the constraints numerically.

\begin{intuition}
Each species is a dye that is carried by the flow, smeared by diffusion at its own rate, and
created or destroyed by chemistry. The equations are coupled only through $\wdot_k(T,\bm Y)$ and
through the flow. With 10 species (\code{h2o2.yaml}) that is 10 equations; with GRI-Mech 3.0, 53.
\end{intuition}

\section{Momentum}

\begin{derivation}[Newton's second law for a control volume]
Momentum per unit volume is $\rho\vu$. The forces are pressure on the surface ($-p\vn\,\dd S$),
viscous stress on the surface ($\bm\tau\cdot\vn\,\dd S$) and gravity in the volume ($\rho\bm g\,\dd V$).
The generic law with $\varphi=\vu$ and Gauss's theorem gives
\begin{equation}
  \pd{(\rho\vu)}{t}+\divg(\rho\vu\vu)=-\nabla p+\divg\bm\tau+\rho\bm g ,
  \label{eq:momentum}
\end{equation}
with $\bm\tau$ the Newtonian stress of \chapref{ch:transport}.
\end{derivation}
\begin{symbols}
$\rho\vu$ & momentum per unit volume & \si{\kilo\gram\per\metre\squared\per\second}\\
$\vu\vu$ & dyadic product: momentum flux tensor per unit density & \si{\metre\squared\per\second\squared}\\
$p$ & pressure & \si{\pascal}\\
$\bm\tau$ & viscous stress tensor & \si{\pascal}\\
$\bm g$ & gravitational acceleration & \si{\metre\per\second\squared}\\
\end{symbols}

\section{Energy}

\begin{derivation}[from total energy to temperature]
\textbf{Step 1: total energy} $e_t=u+\tfrac12|\vu|^2$. Its sources are heat conduction $-\divg\vq$,
work by pressure $-\divg(p\vu)$, work by viscous stresses $\divg(\bm\tau\cdot\vu)$ and gravity
$\rho\bm g\cdot\vu$:
\[
\rho\Dt{e_t}=-\divg\vq-\divg(p\vu)+\divg(\bm\tau\cdot\vu)+\rho\bm g\cdot\vu .
\]
\textbf{Step 2: remove kinetic energy.} Dot \eqref{eq:momentum} with $\vu$ to get
$\rho\,\mathrm D(\tfrac12|\vu|^2)/\mathrm Dt=-\vu\cdot\nabla p+\vu\cdot(\divg\bm\tau)+\rho\bm g\cdot\vu$ and subtract.
Using $h=u+p/\rho$ and mass conservation,
\begin{equation}
  \rho\Dt{h}=\Dt{p}-\divg\vq+\bm\tau:\nabla\vu .
  \label{eq:enthalpy}
\end{equation}
\textbf{Step 3: temperature.} With $h=\sum Y_kh_k(T)$, $\rho\,\mathrm Dh/\mathrm Dt=\rho c_p\,\mathrm DT/\mathrm Dt
+\sum_kh_k\rho\,\mathrm DY_k/\mathrm Dt$. Replace $\rho\,\mathrm DY_k/\mathrm Dt$ by the right-hand
side of \eqref{eq:species} (in the form \eqref{eq:matderiv}). Then expand
$\divg\vq=-\divg(\lambda\nabla T)+\sum_k\divg(h_k\vj_k)$ with $\nabla h_k=c_{p,k}\nabla T$. The
$\sum_kh_k\divg\vj_k$ terms cancel and what remains is
\begin{equation}
  \boxed{\rho c_p\Dt{T}=\Dt{p}+\divg(\lambda\nabla T)-\Big(\sum_k\rho Y_kc_{p,k}\vV_k\Big)\cdot\nabla T
  -\sum_kh_kW_k\wdot_k+\bm\tau:\nabla\vu}
  \label{eq:temperature}
\end{equation}
\end{derivation}
\begin{symbols}
$u$ (scalar) & specific internal energy (not velocity) & \si{\joule\per\kilo\gram}\\
$e_t=u+\tfrac12|\vu|^2$ & specific total energy & \si{\joule\per\kilo\gram}\\
$h=u+p/\rho$ & specific enthalpy & \si{\joule\per\kilo\gram}\\
$\vq$ & heat flux & \si{\watt\per\metre\squared}\\
$p\vu$, $\bm\tau\cdot\vu$ & rate of work by pressure and viscous forces (per area) & \si{\watt\per\metre\squared}\\
$\bm\tau:\nabla\vu$ & viscous dissipation (heating) & \si{\watt\per\metre\cubed}\\
$c_p$, $c_{p,k}$ & specific heat capacity of mixture, of species $k$ & \si{\joule\per\kilo\gram\per\kelvin}\\
$\lambda$ & thermal conductivity & \si{\watt\per\metre\per\kelvin}\\
$\vV_k$, $\vj_k$ & diffusion velocity, diffusive mass flux of $k$ & \si{\metre\per\second}, \si{\kilo\gram\per\metre\squared\per\second}\\
$h_kW_k\wdot_k$ & enthalpy carried by chemical production of $k$ & \si{\watt\per\metre\cubed}\\
$\mathrm D/\mathrm Dt$ & material derivative & unit of argument per \si{\second}\\
\end{symbols}

\begin{intuition}[reading the temperature equation term by term]
\begin{itemize}
  \item $\rho c_p\,\mathrm DT/\mathrm Dt$: heating of a moving fluid parcel.
  \item $\mathrm Dp/\mathrm Dt$: compression heating (important in engines, negligible in open flames).
  \item $\divg(\lambda\nabla T)$: conduction from hot neighbours. In a flame this preheats the fresh gas.
  \item $-\sum\rho Y_kc_{p,k}\vV_k\cdot\nabla T$: species with different heat capacities diffusing
        through a temperature gradient. Small, but kept by Cantera.
  \item $-\sum h_kW_k\wdot_k=\dot q$: \textbf{heat release} by chemistry, the same term as in the 0D reactor
        \eqref{eq:0dT}.
  \item $\bm\tau:\nabla\vu$: viscous heating, negligible at low speed.
\end{itemize}
\end{intuition}

\section{Closing the system}

Unknowns: $\rho$, $\vu$ (3), $p$, $T$, $Y_k$ ($K$): $K+6$. Equations: mass (1), momentum (3),
energy (1), species ($K-1$ independent, plus $\sum Y_k=1$) and the ideal gas law (1): also $K+6$.
The fluxes ($\bm\tau$, $\vq$, $\vV_k$) come from the transport laws, $h_k$ and $c_p$ from the NASA
polynomials, and $\wdot_k$ from the mechanism. \textbf{Cantera supplies all these closures.} The
flow solver only has to discretise derivatives.

\section{The low-Mach-number approximation}
\label{sec:lowmach}

A laminar methane flame moves at \SI{0.4}{\metre\per\second} and the sound speed is
\SI{350}{\metre\per\second}: Mach number $\mathrm{Ma}\sim10^{-3}$. Expanding $p=p_0(t)+\mathrm{Ma}^2p_1(\vx,t)$:
\begin{itemize}
  \item The \textbf{thermodynamic pressure} $p_0$ is spatially uniform and is used in the ideal gas law
        and all property evaluations.
  \item The \textbf{hydrodynamic pressure} $p_1$ (tiny, a few Pa across a flame) appears only in the
        momentum equation to drive the flow.
  \item In an open domain $p_0$ is constant, so $\mathrm Dp/\mathrm Dt\approx0$; viscous heating is
        dropped as well.
\end{itemize}
Sound waves are filtered out, which allows time steps set by the flow rather than by acoustics.
Cantera's 1D flames use exactly this: \code{flame.P} is a single number.

\begin{symbols}
$\mathrm{Ma}=u/c$ & Mach number & \nodim\\
$c$ & speed of sound & \si{\metre\per\second}\\
$p_0(t)$ & thermodynamic (uniform) pressure & \si{\pascal}\\
$p_1(\vx,t)$ & hydrodynamic pressure perturbation & \si{\pascal}\\
\end{symbols}

\section{Why we do not start in 3D}

For \ce{H_2}/air that is already 16 coupled nonlinear PDEs. For methane with GRI-Mech 3.0 it is 59.
The chemical source terms are stiff, flames are thinner than \SI{1}{\milli\metre} while
combustors are metres across, and in turbulent flows the flame is wrinkled at every scale. Before
any of that, two 1D model problems isolate the physics:
\begin{itemize}
  \item \textbf{Premixed flame} (\chapref{ch:premixed}): reactions and diffusion balancing
        convection. Output: flame speed and thickness.
  \item \textbf{Non-premixed counterflow flame} (\chapref{ch:nonpremixed}): mixing controlled by
        strain. Output: flame structure in mixture-fraction space and extinction.
\end{itemize}
Both are also the building blocks (``flamelets'') of turbulent combustion models.
